\chapter{Funkcje specjalne}
\label{ma1:ch:special-functions}

Funkcje specjalne często definiuje się przez całki, których wartości nie dają się wygodnie zapisać za pomocą funkcji elementarnych. Będziemy badać rzeczywiste parametry i wykorzystywać kryteria z rozdziałów \ref{ma1:ch:improper-integrals} i \ref{ma1:ch:parameter-integrals}. Wszystkie całki niewłaściwe rozpatrujemy w sensie Riemanna, niezależnie przy każdym niewłaściwym końcu.

\section{Funkcja Gamma Eulera}
\begin{definicja}
\label{ma1:def:gamma}
Dla $s>0$ funkcją Gamma Eulera nazywamy
\[
 \Gamma(s)=\int_0^\infty x^{s-1}e^{-x}\,\dd x.
\]
\end{definicja}
\begin{twierdzenie}
\label{ma1:thm:gamma-properties}
Całka definiująca $\Gamma(s)$ jest zbieżna i dodatnia dla każdego $s>0$. Zachodzą
\[
 \Gamma(1)=1,\qquad \Gamma(s+1)=s\Gamma(s),\qquad
 \Gamma(n+1)=n!\quad(n\ge0).
\]
\end{twierdzenie}
\begin{proof}
Na $(0,1]$ funkcja podcałkowa jest nieujemna i nie przekracza $x^{s-1}$, którego całka jest skończona, ponieważ $s>0$ (przykład \ref{ma1:ex:improper-p-integrals}). Na $[1,\infty)$ wybierzmy całkowite $N>s+1$. Ze szeregu wykładniczego $e^x\ge x^N/N!$, więc $x^{s-1}e^{-x}\le N!x^{s-1-N}$. Wykładnik jest mniejszy od $-1$, co daje zbieżność przez porównanie. Dodatniość wynika stąd, że na $[1,2]$ ciągła funkcja podcałkowa osiąga ściśle dodatnie minimum, a na pozostałej części dziedziny jest nieujemna. Bezpośrednio $\Gamma(1)=\lim_{R\to\infty}(1-e^{-R})=1$.

Na $[\eta,R]$, $0<\eta<R$, całkowanie przez części daje
\[
 \int_\eta^R x^s e^{-x}\,\dd x
 =-R^s e^{-R}+\eta^s e^{-\eta}
 +s\int_\eta^R x^{s-1}e^{-x}\,\dd x.
\]
Dla $\eta\to0^+$ składnik $\eta^s e^{-\eta}$ dąży do zera. Dla $R\to\infty$ mamy $R^s e^{-R}\le N!R^{s-N}\to0$. Przechodzimy do obu granic niezależnie, korzystając z już udowodnionej zbieżności. Otrzymujemy równanie funkcyjne; indukcja od $n=0$ daje wzór z silnią.
\end{proof}
\begin{twierdzenie}
\label{ma1:thm:gamma-smooth}
Funkcja $\Gamma$ jest klasy $C^\infty$ na $(0,\infty)$, a dla każdej liczby całkowitej $k\ge0$
\[
 \Gamma^{(k)}(s)=\int_0^\infty x^{s-1}(\log x)^ke^{-x}\,\dd x.
\]
\end{twierdzenie}
\begin{proof}
Niech $K=[c,d]\subset(0,\infty)$, $c>0$. Dla $0<x\le1$ i $s\in K$ moduł funkcji podcałkowej po $k$ różniczkowaniach nie przekracza $x^{c-1}|\log x|^k$. Podstawienie $u=-\log x$ na obcięciach daje całkę $\int_0^\infty u^ke^{-cu}\,\dd u=k!/c^{k+1}$ z przykładu \ref{ma1:ex:parameter-exponential}. Dla $x\ge1$ mamy $\log x\le x$: funkcja $x-\log x$ jest niemalejąca na $[1,\infty)$ i przyjmuje wartość $1$ w $1$. Zatem majorantą jest $x^{d+k-1}e^{-x}$, całkowalna z poprzedniego dowodu. Wszystkie pochodne względem $s$ są ciągłe na obciętych przedziałach. Na części $(0,1]$ stosujemy odbicie $x=1-u$, tak że niewłaściwy jest prawy koniec przedziału $[0,1)$, a na $[1,\infty)$ stosujemy bezpośrednio twierdzenia \ref{ma1:thm:uniform-improper-majorant} i \ref{ma1:thm:improper-parameter-differentiation}. Dla każdego rzędu całka z poprzedniej pochodnej jest już zbieżna dzięki odpowiedniej majorancie. Indukcja daje wzór i ciągłość każdej pochodnej.
\end{proof}
W szczególności $\Gamma''(s)>0$, ponieważ funkcja podcałkowa dla $k=2$ jest nieujemna i dodatnia na $[2,3]$. Ponadto $s\Gamma(s)=\Gamma(1+s)\to1$ przy $s\to0^+$, więc $\Gamma(s)\sim1/s$.

\section{Funkcja Beta Eulera}
\begin{definicja}
\label{ma1:def:beta}
Dla $p,q>0$ definiujemy
\[
 B(p,q)=\int_0^1 x^{p-1}(1-x)^{q-1}\,\dd x.
\]
\end{definicja}
Na $(0,1/2]$ czynnik $(1-x)^{q-1}$ jest ograniczony przez $\max(1,2^{1-q})$, a $x^{p-1}$ ma zbieżną całkę, gdyż $p>0$. Na $[1/2,1)$ czynnik $x^{p-1}$ jest ograniczony przez $\max(1,2^{1-p})$, a $(1-x)^{q-1}$ ma zbieżną całkę przez odbicie, gdyż $q>0$. Całka jest więc zbieżna i dodatnia. Podstawienie $x=1-u$ na obcięciach daje $B(p,q)=B(q,p)$.
\begin{twierdzenie}
\label{ma1:thm:beta-recurrence}
Dla $p,q>0$ zachodzą
\[
 B(p,q)=B(p+1,q)+B(p,q+1),\qquad
 B(p,q+1)=\frac{q}{p+q}B(p,q).
\]
\end{twierdzenie}
\begin{proof}
Pierwszy wzór wynika z rozłożenia $1=x+(1-x)$ w funkcji podcałkowej. Pochodna $x^p(1-x)^q$ na $(0,1)$ wynosi
$p x^{p-1}(1-x)^q-q x^p(1-x)^{q-1}$. Całkując na $[\eta,1-\delta]$ i przechodząc niezależnie do granic $\eta,\delta\to0^+$, otrzymujemy
$pB(p,q+1)=qB(p+1,q)$, gdyż funkcja $x^p(1-x)^q$ dąży do zera na obu końcach, a obie całki są zbieżne. Wstawienie $B(p+1,q)=B(p,q)-B(p,q+1)$ daje drugą równość.
\end{proof}

\begin{twierdzenie}[Wzór graniczny Eulera]
\label{ma1:thm:gamma-euler-limit}
Dla $p>0$ mamy
\[
 \Gamma(p)=\lim_{n\to\infty}
 \frac{n^p n!}{p(p+1)\cdots(p+n)}.
\]
\end{twierdzenie}
\begin{proof}
Najpierw wykażemy, że dla ustalonych $p,q>0$
\begin{equation}
 \label{ma1:eq:beta-scaling-limit}
 n^p B(p,q+n)\longrightarrow\Gamma(p).
\end{equation}
Podstawiając $x=u/n$ na obcięciach, dostajemy
$n^p B(p,q+n)=\int_0^n u^{p-1}(1-u/n)^{q+n-1}\,\dd u$.
Dla $n\ge2$ wykładnik $q+n-1$ jest co najmniej $n/2>0$. Rozszerzamy funkcję podcałkową przez zero dla $u\ge n$. Dla $0<u<n$ nierówność $\log(1-v)\le-v$ przy $0\le v<1$ (wynika z pochodnej $\log(1-v)+v$, równej $-v/(1-v)$) daje
\[
 (1-u/n)^{q+n-1}\le e^{-u/2}.
\]
Na każdym $[\eta,R]\subset(0,\infty)$ funkcje dążą jednostajnie do $u^{p-1}e^{-u}$: dla $n>2R$,
\[
 |\log(1-u/n)+u/n|
 =\int_0^{u/n}\frac{v}{1-v}\,\dd v\le (u/n)^2.
\]
Dokładniej, dla $\eta\le u\le R$ błąd spełnia
\[
 |(q+n-1)\log(1-u/n)+u|
 \le |q-1|R/n+(q+n-1)R^2/n^2,
\]
a prawa strona dąży do zera. Stąd $(q+n-1)\log(1-u/n)\to-u$ jednostajnie, a ciągłość jednostajna wykładniczej na ograniczonym przedziale zachowuje tę zbieżność. Wspólna majoranta $u^{p-1}e^{-u/2}$ jest całkowalna na $(0,\infty)$ przez podstawienie $u=2v$ i twierdzenie \ref{ma1:thm:gamma-properties}. Dla $\varepsilon>0$ wybieramy $\eta,R$ tak, aby suma całek majoranty na $(0,\eta)$ i $(R,\infty)$ była mniejsza od $\varepsilon/4$. Wkłady obu ogonów do różnicy całek ciągu i granicy mają wtedy łączny moduł mniejszy od $\varepsilon/2$. Na $[\eta,R]$ twierdzenie \ref{ma1:thm:riemann-uniform-limit} zapewnia dla dostatecznie dużych $n$ moduł różnicy całek mniejszy od $\varepsilon/2$. Nierówność trójkąta dowodzi \eqref{ma1:eq:beta-scaling-limit}.
Dla $q=1$ mamy $B(p,1)=1/p$. Kolejne użycia twierdzenia \ref{ma1:thm:beta-recurrence} dają
$B(p,n+1)=n!/[p(p+1)\cdots(p+n)]$. Podstawienie do \eqref{ma1:eq:beta-scaling-limit} kończy dowód.
\end{proof}
\begin{twierdzenie}
\label{ma1:thm:beta-gamma}
Dla $p,q>0$
\[
 B(p,q)=\frac{\Gamma(p)\Gamma(q)}{\Gamma(p+q)}.
\]
\end{twierdzenie}
\begin{proof}
Z rekurencji Beta dla $n\ge1$
\[
 B(p,q+n)=B(p,q)\prod_{j=0}^{n-1}\frac{q+j}{p+q+j}.
\]
Iloraz dwóch wzorów granicznych Eulera (mianownik dąży do dodatniej liczby) daje
\[
 \frac{\Gamma(p+q)}{\Gamma(q)}
 =\lim_{n\to\infty}n^p\prod_{j=0}^n\frac{q+j}{p+q+j}.
\]
Ostatni czynnik $(q+n)/(p+q+n)$ dąży do $1$, więc pominięcie go nie zmienia granicy. Mnożąc rekurencję przez $n^p$ i stosując \eqref{ma1:eq:beta-scaling-limit}, otrzymujemy
$\Gamma(p)=B(p,q)\Gamma(p+q)/\Gamma(q)$. Dzielimy przez dodatnie wartości Gamma. Dowód nie wymaga zamiany kolejności całkowania w całce podwójnej.
\end{proof}

\section{Całki trygonometryczne i wzór Wallisa}
\begin{twierdzenie}
\label{ma1:thm:beta-trigonometric}
Dla $p,q>0$
\[
 B(p,q)=2\int_0^{\pi/2}\sin^{2p-1}t\,\cos^{2q-1}t\,\dd t.
\]
W szczególności $B(1/2,1/2)=\pi$.
\end{twierdzenie}
\begin{proof}
Podstawienie $x=\sin^2 t$ na domkniętym podprzedziale $(0,\pi/2)$ ma pochodną $2\sin t\cos t$ i prowadzi do podanego wyrażenia. Przy dojściu do końców $t=0,\pi/2$ zmienna $x$ dąży odpowiednio do $0,1$; oba końce rozpatrujemy niezależnie, korzystając ze zbieżności Beta. Dla $p=q=1/2$ funkcja pod całką jest stałą $1$ we wnętrzu, więc całka po pomnożeniu przez $2$ wynosi $\pi$.
\end{proof}
\begin{twierdzenie}[Wzór Wallisa]
\label{ma1:thm:wallis}
\[
 \lim_{n\to\infty}\prod_{k=1}^n
 \frac{(2k)^2}{(2k-1)(2k+1)}=\frac{\pi}{2}.
\]
\end{twierdzenie}
\begin{proof}
Niech $J_m=\int_0^{\pi/2}\sin^m t\,\dd t$ dla całkowitego $m\ge0$. Mamy $J_0=\pi/2$, $J_1=1$. Z redukcji sinusowej (twierdzenie \ref{ma1:thm:sine-reduction}) lub całkowania przez części otrzymujemy dla $m\ge2$:
$J_m=(m-1)J_{m-2}/m$, ponieważ składnik brzegowy $-\sin^{m-1}t\cos t$ znika. Wobec tego dla $n\ge1$
\[
 J_{2n}=\frac{\pi}{2}\prod_{k=1}^n\frac{2k-1}{2k},
 \qquad J_{2n+1}=\prod_{k=1}^n\frac{2k}{2k+1}.
\]
Ponieważ $0\le\sin t\le1$, mamy $J_{2n+1}\le J_{2n}\le J_{2n-1}$. Wszystkie całki są dodatnie. Rekurencja daje $J_{2n-1}/J_{2n+1}=(2n+1)/(2n)$, więc
\[
 1\le J_{2n}/J_{2n+1}\le(2n+1)/(2n).
\]
Iloraz dąży do $1$. Z jawnych iloczynów wynika, że iloczyn w tezie jest równy $(\pi/2)J_{2n+1}/J_{2n}$, co kończy dowód.
\end{proof}

\section{Całka Gaussa i funkcja błędu}
\begin{twierdzenie}
\label{ma1:thm:gaussian-integral}
Dla $a>0$
\[
 \Gamma(1/2)=\sqrt{\pi},\qquad
 \int_{-\infty}^\infty e^{-ax^2}\,\dd x=\sqrt{\pi/a}.
\]
\end{twierdzenie}
\begin{proof}
Z twierdzeń \ref{ma1:thm:beta-gamma} i \ref{ma1:thm:beta-trigonometric} wynika $\Gamma(1/2)^2/\Gamma(1)=\pi$. Ponieważ $\Gamma(1)=1$ i $\Gamma(1/2)>0$, otrzymujemy pierwszy wzór. Podstawienie $u=ax^2$ na $[\eta,R]$, $0<\eta<R$, daje
\[
 \int_\eta^R e^{-ax^2}\,\dd x
 =\frac{1}{2\sqrt a}\int_{a\eta^2}^{aR^2}u^{-1/2}e^{-u}\,\dd u.
\]
Przechodzimy niezależnie do granic $\eta\to0^+$ i $R\to\infty$. Lewa funkcja jest ciągła w zerze, prawa całka Gamma jest zbieżna. Z parzystości i niezależnej zbieżności obu półprostych całka na całej prostej jest dwukrotnością całki od zera. Dodatkowo rekurencja Gamma daje $\Gamma(n+1/2)=\sqrt\pi(2n)!/(4^n n!)$ dla całkowitych $n\ge0$: iloczyn kolejnych czynników $(2k-1)/2$ jest równy $(2n)!/(4^n n!)$ przez wydzielenie czynników parzystych z $(2n)!$.
\end{proof}

\begin{definicja}
\label{ma1:def:erf}
Dla $x\in\R$ definiujemy funkcję błędu i jej dopełnienie:
\[
 \operatorname{erf}(x)=\frac2{\sqrt\pi}\int_0^x e^{-t^2}\,\dd t,
 \qquad \operatorname{erfc}(x)=1-\operatorname{erf}(x).
\]
\end{definicja}
Z podstawowego twierdzenia rachunku całkowego $\operatorname{erf}'(x)=2e^{-x^2}/\sqrt\pi>0$. Funkcja jest nieparzysta (odbicie i parzystość funkcji podcałkowej), a jej granice w $\pm\infty$ to odpowiednio $\pm1$ z całki Gaussa. W szczególności $-1<\operatorname{erf}(x)<1$ dla każdego skończonego $x$. Z addytywności całki
\[
 \operatorname{erfc}(x)=\frac2{\sqrt\pi}\int_x^\infty e^{-t^2}\,\dd t
\]
dla wszystkich $x\in\R$.
\begin{twierdzenie}
\label{ma1:thm:erf-series-tail}
Dla $x\in\R$ zachodzi
\[
 \operatorname{erf}(x)=\frac2{\sqrt\pi}
 \sum_{n=0}^\infty\frac{(-1)^n x^{2n+1}}{n!(2n+1)}.
\]
Dla $x\to+\infty$ mamy
$\operatorname{erfc}(x)\sim e^{-x^2}/(x\sqrt\pi)$.
\end{twierdzenie}
\begin{proof}
Na $[-M,M]$ szereg $\sum (-1)^nt^{2n}/n!$ dąży jednostajnie do $e^{-t^2}$: moduły ogona są ograniczone przez ogon zbieżnego szeregu $\sum M^{2n}/n!$. Całkujemy wyraz po wyrazie na odcinku między $0$ i $x$, stosując twierdzenie \ref{ma1:thm:riemann-uniform-limit} oraz orientację całki dla $x<0$. To daje pierwszy wzór.

Niech $x>0$ i $T(x)=\int_x^\infty e^{-t^2}\,\dd t$. Całkowanie przez części na $[x,R]$, z czynnikami $1/(2t)$ i $2te^{-t^2}$, po przejściu do granicy daje
\[
 T(x)=\frac{e^{-x^2}}{2x}-\int_x^\infty\frac{e^{-t^2}}{2t^2}\,\dd t.
\]
Składnik brzegowy w nieskończoności znika, a ostatnia całka jest zbieżna przez porównanie z $e^{-t^2}/(2x^2)$. Oznaczając $A=e^{-x^2}/(2x)$, mamy $0\le A-T(x)\le T(x)/(2x^2)\le A/(2x^2)$. Stąd $T(x)/A\to1$, co po pomnożeniu przez $2/\sqrt\pi$ daje tezę.
\end{proof}

\begin{center}
\begin{tikzpicture}[x=1.05cm,y=2cm,>=Stealth]
 \draw[->] (-3.2,0)--(3.3,0) node[right] {$x$};
 \draw[->] (0,-1.2)--(0,1.2) node[above] {$y$};
 \draw[dashed,gray] (-3.1,1)--(3.1,1);
 \draw[dashed,gray] (-3.1,-1)--(3.1,-1);
 \draw[blue!70!black,thick,smooth] plot coordinates {(-3,-.99998) (-2.5,-.99959) (-2,-.99532) (-1.5,-.96611) (-1,-.84270) (-.75,-.71116) (-.5,-.52050) (-.25,-.27633) (0,0) (.25,.27633) (.5,.52050) (.75,.71116) (1,.84270) (1.5,.96611) (2,.99532) (2.5,.99959) (3,.99998)};
 \node[above,blue!70!black] at (1.9,.72) {$\operatorname{erf}(x)$};
 \node[left] at (0,1) {$1$};
 \node[right] at (0,-1) {$-1$};
\end{tikzpicture}

\small Rysunek: przybliżony wykres funkcji błędu (punkty zaokrąglone do pięciu miejsc po przecinku), z asymptotami poziomymi.
\end{center}

\section{Całki Fresnela}
\begin{definicja}
\label{ma1:def:fresnel}
Stosujemy normalizację
\[
 C(x)=\int_0^x\cos(t^2)\,\dd t,\qquad
 S(x)=\int_0^x\sin(t^2)\,\dd t.
\]
W innej często używanej normalizacji definiujemy $C_{\pi}(x)=\int_0^x\cos(\pi t^2/2)\,\dd t$ i analogicznie $S_{\pi}$. Podstawienie $u=\sqrt{\pi/2}\,t$ daje
\[
 C_{\pi}(x)=\sqrt{2/\pi}\,C(\sqrt{\pi/2}\,x),\qquad
 S_{\pi}(x)=\sqrt{2/\pi}\,S(\sqrt{\pi/2}\,x).
\]
\end{definicja}
Obie funkcje są nieparzyste, klasy $C^\infty$, a ich pochodne to odpowiednio $\cos(x^2),\sin(x^2)$. Zbieżność przy $x\to+\infty$ wynika z podstawienia $u=t^2$ na $[1,R]$ i kryterium Dirichleta \ref{ma1:thm:improper-dirichlet} dla $\cos u$ lub $\sin u$ oraz $1/(2\sqrt u)$. Funkcje całek częściowych sinusa i cosinusa są ograniczone, a mnożnik maleje do zera.

\begin{twierdzenie}
\label{ma1:thm:fresnel-values}
\[
 \int_0^\infty\cos(t^2)\,\dd t
 =\int_0^\infty\sin(t^2)\,\dd t=\sqrt{\frac{\pi}{8}}.
\]
\end{twierdzenie}
\begin{proof}
Dla ustalonego $a>0$ i $b\in\R$ połóżmy
\[
 Z(b)=\int_0^\infty e^{-at^2}\big(\cos(bt^2)+i\sin(bt^2)\big)\,\dd t.
\]
Całka zespolona oznacza parę całek rzeczywistych. Moduł funkcji podcałkowej jest $e^{-at^2}$, a moduł jej pochodnej względem $b$ jest $t^2e^{-at^2}$. Obie majoranty są całkowalne: pierwsza z twierdzenia \ref{ma1:thm:gaussian-integral}, druga z podstawienia $u=at^2$, które daje $\int_0^\infty t^2e^{-at^2}\,\dd t=\Gamma(3/2)/(2a^{3/2})$. Z twierdzenia \ref{ma1:thm:improper-parameter-differentiation} stosowanego do obu składowych otrzymujemy
$Z'(b)=i\int_0^\infty t^2e^{-at^2}(\cos(bt^2)+i\sin(bt^2))\,\dd t$.
Oznaczmy funkcję podcałkową $Z$ przez $E(t)$. Różniczkując obie składowe, mamy $E'(t)=-2(a-ib)tE(t)$. Całkowanie pochodnej $tE(t)$ na $[0,R]$ daje
\[
 RE(R)=\int_0^R E(t)\,\dd t-2(a-ib)\int_0^R t^2E(t)\,\dd t.
\]
Ponieważ $|RE(R)|=Re^{-aR^2}\to0$ (na przykład $e^{aR^2}\ge aR^2$), po przejściu do granicy otrzymujemy $Z'(b)=iZ(b)/(2(a-ib))$.

Niech $\theta(b)=\arctan(b/a)$ i
\[
 W(b)=\frac{\sqrt\pi}{2}(a^2+b^2)^{-1/4}
 \big(\cos(\theta(b)/2)+i\sin(\theta(b)/2)\big).
\]
Funkcja $W$ jest niezerowa; $\theta'(b)=a/(a^2+b^2)$. Różniczkując czynnik rzeczywisty i trygonometryczny, dostajemy
\[
 \frac{W'(b)}{W(b)}=\frac{-b+ia}{2(a^2+b^2)}
 =\frac{i}{2(a-ib)}.
\]
Ponieważ $Z(0)=\sqrt\pi/(2\sqrt a)=W(0)$ z całki Gaussa, iloraz $Z/W$ ma pochodną zero i wartość $1$ w zerze. Jego obie składowe są stałe z twierdzenia Lagrange'a, więc $Z=W$. Nie korzystamy tutaj z analitycznego przedłużania funkcji Gamma ani z całkowania po krzywych zespolonych.

Pozostaje przejść do granicy $a\to0^+$ dla $b=1$. Po podstawieniu $u=t^2$ na ogonie całkujemy $\cos u$ lub $\sin u$ pomnożone przez $g_a(u)=e^{-au}/(2\sqrt u)$ dla $u\ge1$. Funkcja $g_a$ jest nierosnąca i dąży do zera dla każdego $a\ge0$. Dowód kryterium Dirichleta daje dla $V>U\ge1$
\[
 \left|\int_U^V g_a(u)\cos u\,\dd u\right|
 \le 4g_a(U)\le 2/\sqrt U
\]
i to samo dla sinusa (całki częściowe obu funkcji od $1$ mają moduł co najwyżej $2$). Oszacowanie jest niezależne od $a$. W pierwotnej zmiennej $t$ daje ono moduł całki od $r$ do $s$, $1\le r<s$, nie większy niż $2/r$, ponieważ $U=r^2$. Jednostajne kryterium Cauchy'ego \ref{ma1:thm:uniform-improper-cauchy} daje zbieżność jednostajną całek tłumionych na $a\in[0,\infty)$. Na każdym skończonym przedziale $t\in[0,R]$ mamy $e^{-at^2}\to1$ jednostajnie, bo $0\le1-e^{-at^2}\le aR^2$, z twierdzenia Lagrange'a. Zamiana granicy z całką niewłaściwą \ref{ma1:thm:improper-parameter-limit} pozwala więc przejść do granicy w obu składowych. Dla $b=1$ wzór na $W$ dąży do $\sqrt\pi(\cos(\pi/4)+i\sin(\pi/4))/2$, którego obie składowe wynoszą $\sqrt{\pi/8}$.
\end{proof}
\begin{wniosek}
\label{ma1:cor:fresnel-tail}
Dla $x\to+\infty$, przy $L=\sqrt{\pi/8}$,
\[
 C(x)=L+\frac{\sin(x^2)}{2x}+O(x^{-3}),\qquad
 S(x)=L-\frac{\cos(x^2)}{2x}+O(x^{-3}).
\]
\end{wniosek}
\begin{proof}
Na $[x,R]$ dla $x>0$ używamy $(\sin t^2)'=2t\cos t^2$ i $(\cos t^2)'=-2t\sin t^2$. Całkowanie przez części i przejście do granicy daje
\[
 \int_x^\infty\cos(t^2)\,\dd t=-\frac{\sin(x^2)}{2x}+\int_x^\infty\frac{\sin(t^2)}{2t^2}\,\dd t,
\]
\[
 \int_x^\infty\sin(t^2)\,\dd t=\frac{\cos(x^2)}{2x}-\int_x^\infty\frac{\cos(t^2)}{2t^2}\,\dd t.
\]
Obie całki resztowe są bezwzględnie zbieżne. Jeszcze jedno całkowanie przez części daje
\[
 \int_x^\infty\frac{\sin(t^2)}{2t^2}\,\dd t
 =\frac{\cos(x^2)}{4x^3}
 -\frac34\int_x^\infty\frac{\cos(t^2)}{t^4}\,\dd t,
\]
\[
 \int_x^\infty\frac{\cos(t^2)}{2t^2}\,\dd t
 =-\frac{\sin(x^2)}{4x^3}
 +\frac34\int_x^\infty\frac{\sin(t^2)}{t^4}\,\dd t.
\]
Składniki brzegowe w nieskończoności znikają. Moduł każdej reszty nie przekracza $1/(4x^3)+(3/4)\int_x^\infty t^{-4}\,\dd t=1/(2x^3)$. Odejmując ogony od $L$, otrzymujemy tezę.
\end{proof}

\section{Całka sinusowa i całka cosinusowa}
\begin{definicja}
\label{ma1:def:si-ci}
Dla $x\in\R$ oraz, odpowiednio, $x>0$, definiujemy
\[
 \operatorname{Si}(x)=\int_0^x\frac{\sin t}{t}\,\dd t,
 \qquad
 \operatorname{Ci}(x)=-\int_x^\infty\frac{\cos t}{t}\,\dd t.
\]
W pierwszej całce przyjmujemy wartość funkcji podcałkowej $1$ w $t=0$, zgodnie z granicą $\sin t/t\to1$. Druga całka jest zbieżna z kryterium Dirichleta.
\end{definicja}
Z parzystości $\sin t/t$ funkcja Si jest nieparzysta. Z podstawowego twierdzenia rachunku całkowego $\operatorname{Si}'(x)=\sin x/x$ dla $x\ne0$ i $\operatorname{Si}'(0)=1$. Dla $x,x+h>0$ addytywność daje $\operatorname{Ci}(x+h)-\operatorname{Ci}(x)=\int_x^{x+h}\cos t/t\,\dd t$, więc $\operatorname{Ci}'(x)=\cos x/x$. Z kryterium Dirichleta istnieje także granica Si w $+\infty$.

\begin{twierdzenie}
\label{ma1:thm:sine-integral-value}
\[
 \int_0^\infty\frac{\sin t}{t}\,\dd t=\frac{\pi}{2}.
\]
\end{twierdzenie}
\begin{proof}
Dla $a>0$ i $b\in\R$ definiujemy
$H(b)=\int_0^\infty e^{-at}\sin(bt)/t\,\dd t$, z wartością $b$ w zerze. Funkcja podcałkowa jest ciągła względem pary również w zerze: $\sin(bt)/t=b\,\sin(bt)/(bt)$ z ciągłym przedłużeniem ilorazu przez $1$. Pochodna względem $b$ to $e^{-at}\cos(bt)$, w zerze równa $1$. Jest ciągła i ma całkowalną majorantę $e^{-at}$. Całka w punkcie $b=0$ jest zerowa, więc twierdzenie \ref{ma1:thm:improper-parameter-differentiation} daje
\[
 H'(b)=\int_0^\infty e^{-at}\cos(bt)\,\dd t=\frac{a}{a^2+b^2}.
\]
Ostatnia równość wynika z funkcji pierwotnej
$e^{-at}(-a\cos(bt)+b\sin(bt))/(a^2+b^2)$, której pochodną obliczamy przez regułę iloczynu: w liczniku otrzymujemy $e^{-at}((a^2+b^2)\cos(bt)+(ab-ab)\sin(bt))$; wartość w nieskończoności jest zerowa. Z $H(0)=0$ i pochodnej arctangensa wynika $H(b)=\arctan(b/a)$.

Dla $b=1$ i $a\ge0$ na $[1,\infty)$ mnożnik $g_a(t)=e^{-at}/t$ jest nierosnący i dąży do zera. Kryterium Dirichleta daje dla $V>U\ge1$ oszacowanie $|\int_U^V g_a(t)\sin t\,\dd t|\le4/U$, niezależne od $a$. Z jednostajnego kryterium Cauchy'ego \ref{ma1:thm:uniform-improper-cauchy} wynika zbieżność jednostajna na $a\in[0,\infty)$. Na każdym $[0,R]$ zbieżność funkcji podcałkowych przy $a\to0^+$ jest jednostajna, bo $\sin t/t$ jest ciągła i ograniczona, a $|e^{-at}-1|\le aR$. Z twierdzenia \ref{ma1:thm:improper-parameter-limit} wynika, że granica $H(1)=\arctan(1/a)$ przy $a\to0^+$ jest szukaną całką. Wynosi ona $\pi/2$.
\end{proof}
\begin{twierdzenie}
\label{ma1:thm:si-ci-asymptotics}
Dla $x\to+\infty$ zachodzą
\[
 \operatorname{Si}(x)=\frac\pi2-\frac{\cos x}{x}+O(x^{-2}),
 \qquad \operatorname{Ci}(x)=\frac{\sin x}{x}+O(x^{-2}).
\]
Istnieje też skończona stała $c_0$ taka, że przy $x\to0^+$ mamy $\operatorname{Ci}(x)=\log x+c_0+O(x^2)$.
\end{twierdzenie}
\begin{proof}
Całkowanie przez części na $[x,R]$ z $x>0$ i przejście do granicy daje
\[
 \int_x^\infty\frac{\sin t}{t}\,\dd t
 =\frac{\cos x}{x}-\int_x^\infty\frac{\cos t}{t^2}\,\dd t,
 \quad
 \int_x^\infty\frac{\cos t}{t}\,\dd t
 =-\frac{\sin x}{x}+\int_x^\infty\frac{\sin t}{t^2}\,\dd t.
\]
Ponowne całkowanie przez części daje
\[
 \int_x^\infty\frac{\cos t}{t^2}\,\dd t
 =-\frac{\sin x}{x^2}+2\int_x^\infty\frac{\sin t}{t^3}\,\dd t,
\]
\[
 \int_x^\infty\frac{\sin t}{t^2}\,\dd t
 =\frac{\cos x}{x^2}-2\int_x^\infty\frac{\cos t}{t^3}\,\dd t.
\]
Moduł każdej reszty jest zatem nie większy od $x^{-2}+2\int_x^\infty t^{-3}\,\dd t=2x^{-2}$. Wszystkie składniki brzegowe w nieskończoności dążą do zera, a całki po $t^{-2}$ i $t^{-3}$ są bezwzględnie zbieżne. Wstawienie do definicji i wzoru na pełną całkę sinusową daje asymptotyki w nieskończoności.

Dla $0<x<1$ rozdzielamy
\[
 \operatorname{Ci}(x)=\log x+\operatorname{Ci}(1)
 -\int_x^1\frac{\cos t-1}{t}\,\dd t.
\]
Z $1-\cos t=2\sin^2(t/2)\le t^2/2$ (nierówność $|\sin u|\le|u|$ z twierdzenia Lagrange'a) ostatnia funkcja podcałkowa ma moduł nie większy od $t/2$. Definiujemy
$c_0=\operatorname{Ci}(1)-\int_0^1(\cos t-1)/t\,\dd t$.
Różnica $\operatorname{Ci}(x)-\log x-c_0$ jest całką od $0$ do $x$ z $(\cos t-1)/t$ i ma moduł co najwyżej $x^2/4$. To kończy dowód. Identyfikację stałej $c_0$ ze stałą Eulera--Mascheroniego pozostawiamy poza zakresem tego rozdziału.
\end{proof}

Przed użyciem wartości całki specjalnej należy ustalić jej normalizację, dziedzinę parametrów oraz zbieżność przy wszystkich niewłaściwych końcach. W szczególności całki oscylujące mogą być zbieżne mimo braku zbieżności bezwzględnej; tłumienie wykładnicze nie uprawnia do przejścia do granicy bez osobnej kontroli ogonów.
